\documentclass[10pt,a4paper]{amsart}
\usepackage[margin=19mm]{geometry}
\usepackage{amsmath,amssymb,mathtools}
\usepackage[scr=rsfs,cal=euler]{mathalfa}
\usepackage{xfrac}
\usepackage[unicode,hidelinks,bookmarks=false]{hyperref}
\newcommand{\ie}{i.e.\ }
\newcommand{\cf}{cf.\ }
\newcommand{\resp}{respectively\ }
\newcommand{\Subsection}{\subsection}
\providecommand{\nobreakdash}{\nobreak\hbox{-}}
\setlength{\emergencystretch}{2em}
\multlinegap0pt
\usepackage{bm}
\usepackage{bbm}
\usepackage{dsfont}
\usepackage{xspace}
\usepackage{cancel}
\usepackage{mathtools}
\usepackage{amsthm}
\makeatletter
\@ifundefined{theorem}{\newtheorem{theorem}{Theorem}}{}
\@ifundefined{proposition}{\newtheorem{proposition}[theorem]{Proposition}}{}
\makeatother
\newcommand{\IC}{\operatorname{IC}}
\newcommand{\Ocal}{\mathcal O}
\newcommand{\Pp}{\mathbb P}
\newcommand{\Q}{\mathbb Q}
\title[Gopakumar-Vafa Invariants and Macdonald Formula]{Gopakumar-Vafa Invariants and Macdonald Formula}
\author{Lutian Zhao}
\date{Source: arXiv:2306.05547v2 (CC BY 4.0)}
\begin{document}
\maketitle

\noindent\emph{Source and licence.} Gopakumar-Vafa Invariants and Macdonald Formula. Version arXiv:2306.05547v2 (CC BY 4.0), distributed under the Creative Commons Attribution 4.0 International licence, as recorded in the arXiv licence metadata for this exact version and shown on the abstract page https://arxiv.org/abs/2306.05547v2. Full rights notice: licenses/arxiv-2306.05547-CC-BY-4.0.txt.\par\smallskip

\noindent\textit{Editorial scope.} Counted unit: the Proposition whose source label is \texttt{prop:conic-correction-node} in the article (source0.tex lines 3234--3251, source label \texttt{prop:conic-correction-node}), delivered together with the article's own complete original proof of it (source0.tex lines 3253--3316, one \texttt{proof} environment of 64 lines). No mathematics is written, repaired or re-derived: statement and proof are the article's own text line for line. Required context retained uncounted: none. Every definition, display and label of those lines is retained unchanged. Omitted from this package and not redistributed: 1--3233, 3252--3252, 3317--4580. Nothing inside a retained excerpt is omitted or altered: every retained body is delivered exactly as the article's lines are written, so no omission receipt of any kind accompanies this package. Citations: the retained excerpts carry no citation command at all, so no bibliography entry of the article is transcribed and no reference is redistributed. Reference adaptation: none was needed and none was made; every reference inside the delivered unit resolves inside it, so no unresolved reference or citation is produced. Printed slips kept literal: no printed word, symbol, hypothesis or number of the counted statement or of its proof was altered. Layout and class: the article's own class is \texttt{amsart} with packages including amsmath, amssymb, amsthm, mathtools, geometry, enumitem, hyperref; the wrapper is the lane's pinned \texttt{amsart} editorial wrapper at 10pt on A4, which restates the theorem-like environments the retained text needs and adds the article's own macro definitions that the retained text needs (\texttt{\textbackslash IC}, \texttt{\textbackslash Ocal}, \texttt{\textbackslash Pp}, \texttt{\textbackslash Q}), with extra wrapper packages (bm, bbm, dsfont, xspace, cancel, mathtools). The delivered numbering is the unit's own. Assets: no retained span contains a figure environment, a graphics-inclusion command, a PDF-page inclusion, a table or a TikZ picture; no image file is copied into this package, the package declares no asset dependency and no image file is redistributed with it. Licence: version arXiv:2306.05547v2 of this article is distributed under the Creative Commons Attribution 4.0 International licence, as recorded in the arXiv licence metadata for this exact version and shown on the abstract page https://arxiv.org/abs/2306.05547v2. A full rights notice is delivered at licenses/arxiv-2306.05547-CC-BY-4.0.txt. Characters: every retained excerpt is pure ASCII, so the delivered engine, XeLaTeX with its default Latin Modern text font, reports no missing character.
\begin{proposition}[The incidence singularity]\label{prop:conic-correction-node}
The singular locus of $Y$ is $\mathfrak S\cong\mathfrak B$.  For every
$b\in\mathfrak S$, the analytic germ of the pair $(Y,\mathfrak S)$ at $b$
is
\[
 (Y,\mathfrak S)\cong
 \bigl(\mathbb A^6\times
       \{x_1x_2+x_3x_4=0\},\ \mathbb A^6\times\{0\}\bigr).
\]
Thus, at $b\in\mathfrak S$, the stalk of $\IC_Y$ has one copy of $\Q$ in
degrees $-9$ and $-7$, whereas $\Q_Y[9]$ has only the copy in degree
$-9$.  Moreover,
\[
 \chi(Y)=117,\qquad
 \chi(Y,\Q_Y[9])=-117,\qquad
 \chi(Y,\IC_Y)=-132.
\]
\end{proposition}

\begin{proof}
Let $E$ be the rank-four tautological bundle on $H$.  The incidence space is
the projectivized kernel of the evaluation map
$H^0(\Pp^2,\Ocal(2))\otimes\Ocal_H\to E$.  This map has rank four off
$\mathfrak B$ and rank three on $\mathfrak B$.  At $Z\subset L$, its kernel
is $\ell H^0(\Pp^2,\Ocal(1))$, where $\ell$ is an equation of $L$, and its
cokernel is one-dimensional.  Indeed, a length-four scheme fails to impose
four independent conditions on conics precisely when its degree-two
Hilbert function is at most three.  Hilbert-function growth then forces its
degree-one Hilbert function to be at most two, which is equivalent to
containment in a line.  On a line the rank is exactly three.

The normal derivative of the evaluation map is the multiplication pairing
\[
 H^1(L,\Ocal_L(-3))\otimes H^0(L,\Ocal_L(1))
       \longrightarrow H^1(L,\Ocal_L(-2)).
\]
Indeed, normal deformations of $Z$ away from a fixed line form
$H^0(Z,\Ocal_Z(1))$, while motions of the containing line give the image of
$H^0(L,\Ocal_L(1))$.  The sequence
$0\to\Ocal_L(-3)\to\Ocal_L(1)\to\Ocal_Z(1)\to0$ therefore identifies their
quotient, and hence $N_{\mathfrak B/H,b}$, with
$H^1(L,\Ocal_L(-3))$.  The last factor in the display is the cokernel of
evaluation in degree two.  Serre duality
shows that the pairing is perfect.

The remaining kernel direction,
generated by $\ell^2$, has zero normal derivative.  After analytic changes
of bases, write the final row of the evaluation map as $(a,b,c)$.  The
perfect derivative makes $a,b$ regular parameters normal to
$\mathfrak B$.  The corank-one degeneracy ideal is $(a,b,c)$ and has
underlying locus $\mathfrak B$; since $\mathfrak B$ is smooth of
codimension two, $a,b$ generate its ideal scheme-theoretically and
$c\in(a,b)$.  A final column operation gives the exact row $(u,v,0)$.  On
the affine chart around
$[\ell^2]$ in the exceptional $\Pp^2$, the incidence equation is
$ux+vy=0$.  This is the threefold ordinary double point, with the six
coordinates along $\mathfrak B$ as smooth parameters.  The same normal
calculation applies to every length-four divisor on $L$.

For four distinct points this normal form can be seen directly.  Take
$L=(y=0)$ with affine coordinate $t$ and points $t_i$.  A nearby conic can
be written $q=(y+\ell(t))^2+r(t)$ with $\deg\ell\le1$ and $\deg r\le2$.
Putting $z_i=y_i+\ell(t_i)$ and eliminating the three coefficients of $r$
leaves the single equation
\[
 \sum_{i=1}^4\frac{z_i^2}{\prod_{j\ne i}(t_i-t_j)}=0.
\]
All four coefficients are nonzero, so this is a rank-four quadric.  The
Serre-duality argument above is its coordinate-free extension across
collisions of the four points.

A small resolution of the threefold node has exceptional fibre $\Pp^1$.
The small-map description of its intersection complex gives the two stated
stalk groups.  Finally, $Y\to H$ has fibre $\Pp^1$ off $\mathfrak B$ and
$\Pp^2$ on $\mathfrak B$.  Since
$\chi(H)=51$ and $\chi(\mathfrak B)=3\cdot5=15$,
\[
 \chi(Y)=2(51-15)+3\cdot15=117.
\]
The IC stalk Euler function is $-1$ off $\mathfrak S$ and $-2$ on
$\mathfrak S$.  Integrating this function gives
$-117-15=-132$.
\end{proof}

\end{document}
